% Ad Familiares 16.1 --- headnote
% ad-familiares-16-01, 3 November 50 BC, en route Patrae to Alyzia

\textit{Cicero to Marcus Tullius Tiro, written at sea on
the third day before the Nones of November 50 BC ---
3 November --- between Patrae and Alyzia. Tiro,
Cicero's freedman secretary, had fallen ill at Patrae;
Cicero, pressed to reach Rome and the unresolved business
of his triumph, sailed on without him, and from the
voyage back across the Ionian Sea sends this --- the
first of the ``Tiro left ill at Patrae'' cluster ---
along with the slave Marion to bring word.}

\textit{The salutation is unusual: \textit{Tullius Tironi
suo s. p. d. et Cicero meus et frater et fratris f.} ---
``Tullius to his Tiro, warm greetings, and my Cicero too,
and my brother, and my brother's son'' --- the whole
travelling household putting its name to the letter, as if
to reinforce by sheer family weight that Tiro is missed.
The body is fatherly anxiety in three short movements:
permission to stay (\textit{approbavi tuum consilium}),
the careful offer of a route to catch up (Marion has been
dispatched to ferry him onward if he is fit, or come back
alone if not), and the close that resolves the apparent
contradiction --- \textit{amor ut valentem videamus
hortatur, desiderium ut quam primum; illud igitur
potius}: love wants him well, longing wants him quickly,
and love wins. The whole cluster is the personal
correspondence of the corpus at its most exposed.}
